% introduction.tex — Introduction for NF-Langevin EPR paper

\section{Introduction}
\label{sec:introduction}

Stochastic thermodynamics provides a unified framework for quantifying the irreversibility of mesoscopic systems driven out of equilibrium~\cite{Seifert2012}.
At its core lies the entropy production rate (EPR), which measures the rate of irreversible heat dissipation and serves as the fundamental diagnostic of detailed-balance violation.
For overdamped Langevin systems with known dynamics, the EPR can be computed exactly from the Fokker--Planck probability current, provided the time-dependent probability density $p(x,t)$ is available~\cite{Sekimoto2010}.
In the context of periodically driven bistable systems---a paradigm for molecular motors, genetic switches, and phase transitions---the EPR encodes the thermodynamic cost of operating the system cyclically and constrains the efficiency of information processing through thermodynamic uncertainty relations~\cite{Barato2015,Gingrich2016}.

A persistent challenge in applying these ideas is that the full probability density is rarely accessible.
When the continuous state space is projected onto a discrete set of mesostates---for instance, by coarse-graining a double-well potential into left and right basins---the resulting master-equation description systematically underestimates the entropy production.
Schnakenberg's network-theoretic formula~\cite{Schnakenberg1976} applied to such a two-state projection captures only the interwell transition entropy, discarding all intrawell probability currents.
As Esposito~\cite{Esposito2012} demonstrated rigorously, this coarse-graining renders a fraction of the dissipation ``hidden,'' and the gap between the total (Sekimoto) and coarse-grained (Schnakenberg) entropy production can be substantial---exceeding \num{90}\% in strongly driven bistable systems.

The problem of recovering hidden entropy production from limited observations has motivated machine-learning approaches that exploit trajectory statistics.
Neural network classifiers trained to distinguish forward from time-reversed trajectories (NEEP~\cite{Kim2020}) provide lower bounds on the entropy production, while score-matching methods have recently been used to learn the probability current directly in steady-state active matter~\cite{Boffi2024}.
However, these approaches are either limited to scalar bounds on the integrated EP or to stationary systems.

Normalizing flows~\cite{Rezende2015,Kobyzev2021,Papamakarios2021}---bijective neural-network maps that transform a simple base distribution into a complex target via the change-of-variables formula---offer an alternative: they provide explicit, differentiable density estimates with tractable likelihoods, enabling direct computation of the score function $\partial\log p/\partial x$ needed to reconstruct the Fokker--Planck current.
In the broader statistical mechanics context, normalizing flows have been deployed as Boltzmann generators for equilibrium sampling~\cite{Noe2019}, for free-energy estimation~\cite{Wirnsberger2020}, and for lattice field theory~\cite{Albergo2019}, but their application to nonequilibrium entropy production estimation remains unexplored.

In this work, we apply one-dimensional planar normalizing flows~\cite{Rezende2015} with a Gaussian mixture base distribution to learn the non-equilibrium density $p(x,t)$ from ensemble trajectory data generated by overdamped Langevin dynamics in a periodically tilted Landau potential $V(x,h) = x^4 - 2x^2 + hx$.
The flow-based density is used to reconstruct the spatially resolved Fokker--Planck probability current and compute the full EPR density $\sigma(x,t) = J^2/(D\,p)$, without any coarse-graining.
We demonstrate four principal results:
\begin{enumerate}
    \item The NF-based EPR recovers the full Sekimoto dissipation ($\Sigma_{\mathrm{NF}} \approx \Sigma_{\mathrm{Sek}}$), confirming that the learned density captures the microscopic probability currents responsible for the vast majority of entropy production.
    \item The two-state Schnakenberg estimate captures only $\approx\num{5}\%$ of the total dissipation, with $\sim\!\num{95}\%$ constituting hidden intrawell entropy production that is recovered by the flow-based method.
    \item A Gaussian mixture base distribution is a necessary precondition for learning bimodal densities with monotone planar flows: by the monotone composition theorem, no number of monotone layers can transform a unimodal base into a multimodal density.
    \item A force-free extension computes the probability current from the continuity equation $\partial p/\partial t = -\partial J/\partial x$ using only the NF-learned densities and the diffusion coefficient~$D$ (estimated via Kramers--Moyal), achieving $\num{11.3}\%$ accuracy without knowledge of the force field.
\end{enumerate}

The entire implementation uses pure NumPy and SciPy, without automatic differentiation or deep\-learning frameworks, making it accessible for pedagogical purposes and small-scale research applications.
The method complements a companion paper on topological hysteresis bounds in the same Landau system~\cite{FernandezCaballero2026}, which established the thermodynamic hierarchy between deterministic and stochastic dissipation using Sekimoto and Schnakenberg methods but did not resolve the intrawell entropy production.

The remainder of this paper is organized as follows.
\Cref{sec:methods} presents the physical model, the three EPR estimation methods, the normalizing-flow architecture, and a force-free extension via the continuity equation.
\Cref{sec:results} validates the approach on equilibrium and harmonic benchmarks before presenting the main EPR comparison and the force-free validation.
\Cref{sec:discussion} discusses the results, limitations, and future directions.

%-----------------------------------------------------------------------------
\subsection*{What is entropy production, physically?}
%-----------------------------------------------------------------------------

For readers unfamiliar with stochastic thermodynamics, we offer an intuitive picture.
\textit{Entropy production} quantifies irreversibility---the thermodynamic ``cost'' of driving a system away from equilibrium.
An everyday analogy is stirring coffee: the work done by the spoon dissipates as heat into the liquid, never spontaneously reversing.
At the molecular scale, this dissipation manifests as probability currents that violate detailed balance.

In a bistable potential like \cref{eq:landau_potential}, driving the system cyclically creates two kinds of dissipation:
\begin{enumerate}
    \item \textbf{Interwell transitions:} particles hopping between the left and right wells---the macroscopic switching events visible in coarse-grained descriptions.
    \item \textbf{Intrawell relaxation:} probability currents \emph{within} each well as the density continuously adjusts to the tilting force---invisible to a two-state projection.
\end{enumerate}
The Schnakenberg formula captures only the first; the second constitutes the ``hidden'' entropy production.
In strongly driven bistable systems, this hidden fraction can exceed \num{95}\% of the total dissipation.
The normalizing-flow method developed here recovers the full spatially-resolved probability current, making the hidden dissipation directly observable.
